% !TeX root = ../教材题库_全部条目_转写工作稿.tex
% Source: XJ-B1-C04-EX-P117-Y1349
% 第4章 / 4.3 对数函数 / 教材 p.117、118 / PDF 第124、125页
% OCR未校对：公式与文字均可能有误；未完成数学TeX转写；禁止直接发给学生
% 原题及相邻原解截图（核对用）：../题目原图/XJ-B1-C04-EX-P117-Y1349-01.webp, ../题目原图/XJ-B1-C04-EX-P117-Y1349-02.webp
% 本文件只在转写完成后作为题干使用；OCR 条目尚不是可靠题干。
\begin{bankitem}{4.3 例题 例题3}{第4章 / 4.3 对数函数 / 教材 p.117、118 / PDF 第124、125页}{XJ-B1-C04-EX-P117-Y1349}{OCR未校对：公式与文字均可能有误；未完成数学TeX转写；禁止直接发给学生}
\begin{ocrdraft}
[教材p.117 / PDF第124页]
\par
0二0：]设A=log\ocrglyph{201E}x,B=log\ocrglyph{201E}y，C=log\ocrglyph{201E}z，
\par
用A,B，C表示下列各式：
\par
（1）logz之3；
\par
（2）
\par
解（Dloga．一一一log\ocrglyph{201E}(xy2)一
\par
3
\par
log\ocrglyph{201E}x十210gzy\_310g\ocrglyph{201E}z
\par
一A十2B一3C；
\par
[教材p.118 / PDF第125页]
\par
==310gux十—lwy\ocrglyph{00AF}\ocrglyph{00AF}lwz
\par
==3A十—B——C\ensuremath{\cdot}
\end{ocrdraft}
\end{bankitem}
